\documentclass{article}


% if you need to pass options to natbib, use, e.g.:
%     \PassOptionsToPackage{numbers, compress}{natbib}
\PassOptionsToPackage{sort, numbers, compress}{natbib}
% before loading neurips_2022

% ready for submission
% \usepackage{neurips_2023}


% to compile a preprint version, e.g., for submission to arXiv, add add the
% [preprint] option:
     \usepackage[preprint]{neurips_2023}


% to compile a camera-ready version, add the [final] option, e.g.:
     %\usepackage[final]{neurips_2023}


% to avoid loading the natbib package, add option nonatbib:
   %\usepackage[nonatbib]{neurips_2023}


\usepackage[utf8]{inputenc} % allow utf-8 input
\usepackage[T1]{fontenc}    % use 8-bit T1 fonts
\usepackage{url}            % simple URL typesetting
\usepackage{booktabs}       % professional-quality tables
\usepackage{amsfonts}       % blackboard math symbols
\usepackage{nicefrac}       % compact symbols for 1/2, etc.
\usepackage{microtype}      % microtypography
\usepackage{xcolor}         % colors


\usepackage{makeidx}
\usepackage{amsmath,amsfonts,amssymb,amsthm}
\usepackage{graphics}
\usepackage{graphicx}
\allowdisplaybreaks
\thispagestyle{empty}
\usepackage{framed}

\usepackage{subfigure}

\usepackage{framed}
\usepackage{mathrsfs}
\usepackage{enumerate}
\usepackage{tikz}
\usetikzlibrary{matrix}
\usepackage[all]{xy}
\usetikzlibrary{shapes}
\usetikzlibrary{decorations.markings}
\usepackage{multirow}
\usepackage{bm}

\usepackage{thmtools}
\usepackage{mathptmx} %font
\usepackage{pgfplots}
\pgfplotsset{compat = newest}

% --- Custom Math Commands --- %
\newtheorem{theorem}{Theorem}
\declaretheoremstyle{lemma}
\declaretheorem[style=lemma, name=Lemma]{lemma}

\theoremstyle{definition}
\newtheorem{definition}{Definition}

\declaretheoremstyle{example}
\declaretheorem[style=example, name=Example]{example}

\theoremstyle{remark}
\newtheorem*{remark}{Remark}

\declaretheoremstyle{proposition}
\declaretheorem[style=proposition, name=Proposition]{proposition}

\declaretheorem[name=Note]{note}
\declaretheoremstyle{note}

\newenvironment{ftheo}
  {\begin{mdframed}\begin{theorem}}
  {\end{theorem}\end{mdframed}}


\definecolor{mydarkblue}{rgb}{0,0.08,0.45}
\usepackage[colorlinks=true,
    linkcolor=mydarkblue,
    citecolor=mydarkblue,
    filecolor=mydarkblue,
    urlcolor=mydarkblue,
    pdfview=FitH]{hyperref}

\usepackage[T1]{fontenc}

\pagestyle{empty}
% ---- Algorithms ----- %
\usepackage[algoruled, linesnumbered, shortend, lined]{algorithm2e}
\usepackage{algpseudocode}
\SetAlgoSkip{bigskip}
\SetKwInput{Input}{Input\hspace{34pt}{ }}
\SetKwInput{Output}{Output\hspace{27pt}{ }}
\SetKwInput{Initialize}{Initialize\hspace{20pt}{ }}
\SetKwInput{Parameters}{Parameters\hspace{2pt}{ }}

\title{Final Project Report MATH 563}


% The \author macro works with any number of authors. There are two commands
% used to separate the names and addresses of multiple authors: \And and \AND.
%
% Using \And between authors leaves it to LaTeX to determine where to break the
% lines. Using \AND forces a line break at that point. So, if LaTeX puts 3 of 4
% authors names on the first line, and the last on the second line, try using
% \AND instead of \And before the third author name.


\author{%
  Alexandre St-Aubin \\
  Student ID: 261115785
  % examples of more authors
  \And
 Jonathan Campana \\
 Student ID: \\
 \AND
 Edmund Wu\\
 Student ID: \\
 \And
 Elliot Forcier-Poirier\\
 Student ID: \\ 
}

% import macros.tex
\input{macros}

\begin{document}

\maketitle


\begin{abstract}
  give some background
\end{abstract}


\section{Introduction}
\subsection{Notation}%
  \label{sub:Notation}
\subsection{Preleminaries}%
  \label{sub:Preleminaries}
  We first note an important relation between the proximal and subdifferential operators.\\
  \begin{proposition}
    Given a function $f\in\Gamma$, we have 
    \begin{equation}
      \prox_{\lambda f} =(I+ \lambda \partial f)^{-1},
    \end{equation}  
    where $(I + \lambda \partial f)^{-1}$, for $\lambda > 0$, is called the \textbf{resolvent} of the operator $\partial f$. In particular, the proximal operator is the resolvent of the subdifferential operator \cite{parikh_2014}.\\
  \end{proposition}
  \begin{definition}[Maximal Monotone operator]
    A monotone operator is \textbf{maximal} if its graph is not properly contained in the graph of any other monotone operator. If A is monotone (resp. maximal monotone) then so are $A^{-1}$ and $\lambda A$ if $\lambda > 0$.
  \end{definition}
  The next proposition, given in \cite{ConnorVander}, will be useful in the construction of the optimizing algorithms.\\ 
  \begin{proposition}     
If $F$ is maximal monotone, then $(I + \lambda F)^{-1}(\hat x)$ exists and is unique for all $\hat x \in R^n$. The value $x = (I + \lambda F)^{-1}(\hat x)$ of the resolvent is the unique solution of the monotone inclusion
 \begin{equation*}
      \hat x \in x + \lambda {F}(x) 
   \end{equation*}
  \end{proposition}  

\section{Algorithms}
The algorithms we present are used to solve the non-blind image deblurring problem, which is a convex optimization problem of the form 
\begin{equation}
  \begin{split}
    \min_{x,y} f(x) + g(y)\\ 
    \text{subject to }Ax= y,
  \end{split}
\end{equation}
where $f, g \in \Gamma$ have inexpensive prox-operators. We are given a vectorized, blurred image $b$, along with the kernel $K$ used to blur it using the model 
\begin{equation}
  Kx + \eta = b,
\end{equation}
where $x$ is the original image, and $\eta$ the added noise. Our goal is to recover the original image $x$. We formulate the problem in two ways, namely, one which uses the $L^1$ norm, and the other the $L^2$ norm. First,
\begin{equation}\label{l1}
    \min_x \|Kx - b\|_1 + \delta_S(x) + \gamma \|Dx \|_{\iso}, 
\end{equation}
and second,
\begin{equation}\label{l2}
    \min_x \|Kx - b\|_2^2 + \delta_S(x) + \gamma \|Dx \|_{\iso}, 
\end{equation}
where $S = \{x \mid 0 \leq x \leq 1\}$, so that $\delta_S$ aims to restrict the range of allowed pixels. The $\iso$ norm is defined as 
$$\|(x,y) \|_{\iso} = \sum^{n}_{k=1} (x_k^2 + y^2_k )^{\frac{1}{2}} $$

\subsection{Primal Douglas-Rachford splitting (Spingarn's method)}%
  \label{sub:Primal Douglas-Rachford splitting (Spingarn's method)}
  This first algorithm, which we will refer to by PDRS, was invented in the 1950s as a method to solve the heat equation \cite{douglas_rachford_1956}. Per (\ref{l1}), let 
  \begin{equation}\label{gf}
    \begin{split}
      f(x) =& \delta_S(x)\\ 
      g(x) =& \|Kx - b\|_1 + \gamma \|Dx\|_\iso
    \end{split}
  \end{equation}
  and consider the minimization problem 
  \begin{equation}
    \begin{split}
      \min_{x,y} f(x) + g(y) +\delta_0 \left( \begin{bmatrix}
        A & -I
      \end{bmatrix} \begin{bmatrix}
        x\\y
      \end{bmatrix} \right)
    \end{split}
  \end{equation}
  with the optimality condition 
  \begin{equation}
    \begin{split}
      0 \in \partial f(x) + \partial g(y) + \partial \delta_0 \left( \begin{bmatrix}
        A & -I
      \end{bmatrix} \begin{bmatrix}
        x\\y
      \end{bmatrix} \right) =: F(x,y) 
    \end{split}
  \end{equation}
  Since this particular problem is too computationally expensive, the strategy is to split $F$ into $F= A+B$. Define 
  \begin{equation}
    \begin{split}
      A(x,y)& = \partial f(x) + \partial g(y) = \partial (\delta_S (x)) +\partial (\|Ky - b\|_1 + \gamma \|Dy\|_{\iso})\\ 
      B(x,y) &= \partial \delta_0 \left( \begin{bmatrix}
        A & -I
      \end{bmatrix} \begin{bmatrix}
        x\\y
      \end{bmatrix} \right)
    \end{split}
  \end{equation}
\newpage
\section*{Acknowledgments}
Please include a list of each person (with student ID) in your group
and write a paragraph \textbf{for each person} with exactly what they
did. Note that the Acknowledgment section does not count towards the
page limit.

Example Acknowledgments would look like:

\paragraph{Alexandre St-Aubin} (Student ID: 261115785) We ball.
\paragraph{Jonathan Campana} (Student ID: ) We ball. 
\paragraph{Elliot Forcier-Poirier} (Student ID: ) We ball. 
\paragraph{Edmund Wu} (Student ID: ) We ball. 
\bibliographystyle{plainnat}
\bibliography{563_FinalReport.bib}

\medskip



%%%%%%%%%%%%%%%%%%%%%%%%%%%%%%%%%%%%%%%%%%%%%%%%%%%%%%%%%%%%


\appendix


\section{Appendix}

Optionally include extra information (complete proofs, additional experiments and plots) in the appendix.
This section will often be part of the supplemental material.


\end{document}

